\documentclass[10pt,letterpaper]{article}
\usepackage{fontspec}
\usepackage{amsmath,amssymb}
\usepackage{graphicx}
\usepackage{xcolor}
\usepackage{tcolorbox}
\tcbuselibrary{skins}
\usepackage{geometry}
\usepackage{booktabs}
\usepackage{colortbl}
\usepackage{fancyhdr}
\usepackage{titlesec}
\usepackage[shortlabels,inline]{enumitem}
\usepackage{tabularx}
\usepackage{array}
\usepackage{multicol}
\usepackage{tikz}

\setmainfont{Times New Roman}
\setsansfont{Arial}

\geometry{
    letterpaper,
    top=1.2cm,
    bottom=1.2cm,
    left=1.2cm,
    right=1.2cm,
    headheight=14pt,
    headsep=0.4cm,
    footskip=0.4cm
}

\definecolor{stewartcyan}{RGB}{0, 118, 186}
\definecolor{stewartred}{RGB}{218, 41, 28}
\definecolor{stewarttableheader}{RGB}{210, 224, 238}
\definecolor{stewarttablehead}{RGB}{210, 224, 238}
\definecolor{tableblue}{RGB}{210, 224, 238}
\definecolor{stewartlightblue}{RGB}{235, 243, 250}
\definecolor{stewartblue}{RGB}{0, 118, 186}
\definecolor{stewartgray}{RGB}{120, 120, 120}
\definecolor{stewartlightgray}{RGB}{240, 240, 240}
\definecolor{stewartpurple}{RGB}{85, 30, 95}
\definecolor{stewartpurplebg}{RGB}{243, 238, 245}
\definecolor{stewartpurpletext}{RGB}{85, 30, 95}
\definecolor{darkgray}{RGB}{40, 40, 40}

\pagestyle{fancy}
\fancyhf{}
\renewcommand{\headrulewidth}{0pt}

\providecommand{\makecell}[2][c]{\begin{tabular}[#1]{@{}c@{}}#2\end{tabular}}
\providecommand{\faDesktop}{\textbf{[Máy tính]}}
\providecommand{\casicon}{\textbf{\textsf{[CAS]}}}
\providecommand{\ticon}{\textbf{\textsf{[T]}}}
\providecommand{\captionof}[2]{\par\vspace{2pt}{\small\textbf{#1}: #2}\par}
\NewDocumentEnvironment{tasks}{o d()}{\begin{enumerate}}{\end{enumerate}}
\providecommand{\task}{\item}
\newenvironment{redframebox}{\begin{definitionbox}}{\end{definitionbox}}

\newtcolorbox{definitionbox}{
    colback=white,
    colframe=stewartred,
    arc=0mm,
    boxrule=0.9pt,
    left=3.5mm, right=3.5mm, top=3mm, bottom=3mm
}

\begin{document}
\documentclass[10pt]{article}
\usepackage[paperwidth=210mm, paperheight=297mm, left=16mm, right=16mm, top=18mm, bottom=16mm]{geometry}
\usepackage{amsmath,amssymb,amsfonts}
\usepackage{fontspec}
\usepackage{xcolor}
\usepackage{paracol}
\usepackage{tcolorbox}
\usepackage{tikz}
\usepackage{enumitem}
\usepackage{fancyhdr}

\IfFontExistsTF{TeX Gyre Termes}{
  \setmainfont{TeX Gyre Termes}
}{
  \IfFontExistsTF{Times New Roman}{\setmainfont{Times New Roman}}{}
}
\IfFontExistsTF{TeX Gyre Heros}{
  \setsansfont{TeX Gyre Heros}
}{
  \IfFontExistsTF{Arial}{\setsansfont{Arial}}{}
}

\definecolor{stewartblue}{RGB}{0, 130, 160}
\definecolor{stewartred}{RGB}{205, 35, 35}

\pagestyle{fancy}
\fancyhf{}
\fancyhead[L]{\textsf{\small 220}\qquad\textsf{\small\textbf{CHƯƠNG 3}}\quad\textsf{\small Các quy tắc tính đạo hàm}}
\renewcommand{\headrulewidth}{0pt}
\setlength{\headheight}{14pt}

\fancyfoot[C]{%
  \fontsize{5.5pt}{6.5pt}\selectfont\color{gray!80!black}%
  \begin{tabular}{@{}c@{}}
  Copyright 2021 Cengage Learning. All Rights Reserved. May not be copied, scanned, or duplicated, in whole or in part. WCN 02-200-203\\[-1pt]
  Copyright 2021 Cengage Learning. All Rights Reserved. May not be copied, scanned, or duplicated, in whole or in part. Due to electronic rights, some third party content may be suppressed from the eBook and/or eChapter(s).\\[-1pt]
  Editorial review has deemed that any suppressed content does not materially affect the overall learning experience. Cengage Learning reserves the right to remove additional content at any time if subsequent rights restrictions require it.
  \end{tabular}%
}

\setlength{\parindent}{0pt}
\setlength{\parskip}{4.5pt plus 1pt minus 1pt}

\begin{document}

\setcolumnwidth{46mm, 122mm}
\setlength{\columnsep}{10mm}

\begin{paracol}{2}
\switchcolumn[1]

Kết quả của Ví dụ 6 rất đáng ghi nhớ:

\vspace{0.2ex}
\noindent
\begin{minipage}{\linewidth}
  \makebox[0pt][l]{%
    \tikz[baseline=-0.6ex]{%
      \node[draw=stewartred, line width=0.8pt, inner sep=2.5pt, minimum size=5.2mm] {\bfseries\textcolor{stewartred}{4}};
    }%
  }%
  \centering
  \tikz[baseline=-0.6ex]{%
    \node[draw=stewartred, line width=0.8pt, inner sep=6pt] {%
      $\displaystyle \frac{d}{dx} \ln|x| = \frac{1}{x}$%
    };
  }%
\end{minipage}
\vspace{0.8ex}

\noindent\textcolor{stewartblue}{\rule{1.2ex}{1.2ex}}\hspace{0.5em}\textbf{\large Đạo hàm logarit}

Việc tính đạo hàm của các hàm số phức tạp chứa tích, thương hoặc lũy thừa thường có thể được đơn giản hóa bằng cách lấy logarit. Phương pháp được sử dụng trong ví dụ sau đây được gọi là \textbf{đạo hàm logarit}.

\vspace{0.5ex}
\noindent\textbf{\textcolor{stewartblue}{VÍ DỤ 7}}\hspace{1em}Tính đạo hàm của $y = \dfrac{x^{3/4}\sqrt{x^2+1}}{(3x+2)^5}$.

\vspace{0.5ex}
\noindent\textbf{\textcolor{stewartblue}{LỜI GIẢI}}\hspace{1em}Lấy logarit tự nhiên ở hai vế của phương trình và sử dụng các Quy tắc của Logarit để đơn giản hóa:
\[
\ln y = \frac{3}{4}\ln x + \frac{1}{2}\ln(x^2+1) - 5\ln(3x+2)
\]
Lấy đạo hàm ẩn theo biến $x$, ta được
\[
\frac{1}{y}\frac{dy}{dx} = \frac{3}{4}\cdot\frac{1}{x} + \frac{1}{2}\cdot\frac{2x}{x^2+1} - 5\cdot\frac{3}{3x+2}
\]
Giải phương trình theo $dy/dx$, ta nhận được
\[
\frac{dy}{dx} = y \left( \frac{3}{4x} + \frac{x}{x^2+1} - \frac{15}{3x+2} \right)
\]

\switchcolumn*
\switchcolumn[0]
\begin{flushleft}
\sffamily\footnotesize\color{black!90}
Nếu chúng ta không sử dụng đạo hàm logarit trong Ví dụ~7, ta sẽ phải dùng cả Quy tắc thương lẫn Quy tắc tích. Phép tính thu được khi đó sẽ cực kỳ cồng kềnh.
\end{flushleft}

\switchcolumn[1]
\noindent Vì đã có một biểu thức tường minh cho $y$, ta có thể thay thế vào và viết
\begin{equation*}
\frac{dy}{dx} = \frac{x^{3/4}\sqrt{x^2+1}}{(3x+2)^5} \left( \frac{3}{4x} + \frac{x}{x^2+1} - \frac{15}{3x+2} \right) \tag*{\textcolor{stewartblue}{\blacksquare}}
\end{equation*}

\vspace{0.3ex}
\begin{tcolorbox}[
  colback=white,
  colframe=stewartred,
  boxrule=0.7pt,
  sharp corners,
  top=6pt, bottom=6pt, left=8pt, right=8pt
]
{\bfseries\textcolor{stewartred}{Các bước trong Đạo hàm logarit}}\par\vspace{3pt}
\begin{enumerate}[label=\textbf{\arabic*.}, leftmargin=1.5em, itemsep=2pt, parsep=0pt, topsep=2pt]
  \item Lấy logarit tự nhiên ở hai vế của phương trình $y = f(x)$ và sử dụng các Quy tắc của Logarit để khai triển biểu thức.
  \item Lấy đạo hàm ẩn theo biến $x$.
  \item Giải phương trình kết quả để tìm $y'$ và thay thế $y$ bởi $f(x)$.
\end{enumerate}
\end{tcolorbox}

Nếu $f(x) < 0$ với một số giá trị của $x$, thì $\ln f(x)$ không xác định, nhưng ta vẫn có thể sử dụng phương pháp đạo hàm logarit bằng cách trước hết viết $|y| = |f(x)|$ rồi sau đó sử dụng Công thức~4. Ta minh họa quy trình này bằng cách chứng minh phiên bản tổng quát của Quy tắc lũy thừa, như đã hứa ở Mục~3.1. Nhắc lại rằng phiên bản tổng quát của Quy tắc lũy thừa khẳng định rằng nếu $n$ là một số thực bất kỳ và $f(x) = x^n$, thì $f'(x) = nx^{n-1}$.

\vspace{0.5ex}
\noindent\textbf{\textcolor{stewartred}{CHỨNG MINH QUY TẮC LŨY THỪA (DẠNG TỔNG QUÁT)}}\hspace{1em}Đặt $y = x^n$ và sử dụng đạo hàm logarit:

\switchcolumn*
\switchcolumn[0]
\vspace{1ex}
\begin{flushleft}
\sffamily\footnotesize\color{black!90}
Nếu $x = 0$, ta có thể chỉ ra rằng $f'(0) = 0$ với $n > 1$ trực tiếp từ định nghĩa của đạo hàm.
\end{flushleft}

\switchcolumn[1]
\[
\ln |y| = \ln |x|^n = n \ln |x| \qquad x \neq 0
\]
\vspace{0.5ex}
\noindent
\begin{tabular*}{\linewidth}{@{}p{2.2cm}@{}l@{\extracolsep{\fill}}r@{}}
Do đó & $\displaystyle \frac{y'}{y} = \frac{n}{x}$ & \\[2.5ex]
Vì vậy & $\displaystyle y' = n \frac{y}{x} = n \frac{x^n}{x} = nx^{n-1}$ & \textcolor{stewartred}{\blacksquare}
\end{tabular*}

\end{paracol}

\end{document}
\end{document}
